% !TeX program = xelatex
% UTF-8. Compile with XeLaTeX from this directory.
% For inclusion in a XePersian document, define \IntussusceptionInNotebook.
% Source: Google Drive Notebook/intussusception; original filename order.
% pic01.jpg = 20261012_201605657.jpg
% pic02.jpg = 20261012_201612173.jpg
% pic03.jpg = 20261012_201619461.jpg
% pic04.jpg = 20261012_201625948.jpg
% pic05.jpg = 20261012_201633240.jpg
% pic06.jpg = 20261012_201644447.jpg
% pic07.jpg = 20261012_201651423.jpg
\ifdefined\IntussusceptionInNotebook
  \let\IntussusceptionFinish\relax
\else
  \def\IntussusceptionFinish{\end{document}}
  \documentclass[12pt,a4paper]{article}
  \usepackage[margin=22mm]{geometry}
  \usepackage{amsmath,amssymb}
  \usepackage{graphicx}
  \usepackage[normalem]{ulem}
  \usepackage{xepersian}
  \IfFontExistsTF{Amiri}{\settextfont{Amiri}}{\settextfont{DejaVu Sans}}
  \setlatintextfont{DejaVu Serif}
  \setlength{\parindent}{0pt}
  \setlength{\parskip}{5pt}
  \setlength{\emergencystretch}{3em}
  \begin{document}
\fi
\providecommand{\IntussusceptionUnclear}[1]{\textbf{[خوانش نامطمئن: #1]}}
\providecommand{\IntussusceptionPage}[2]{\clearpage\section*{تصویر #1}\lr{\texttt{#2}}\par}
\providecommand{\IntussusceptionArrow}{\lr{\(\longleftarrow\)}}
\providecommand{\IntussusceptionImagePath}{images/}
\IfFileExists{images/pic01.jpg}{}{%
  \renewcommand{\IntussusceptionImagePath}{pediatrics/ped-gastroenterology/intussusception/images/}}

\begin{center}
{\Large\bfseries اینتوساسپشن در کودکان}\par
\lr{Intussusception}
\end{center}
\textbf{یادداشت رونویسی:} عبارت‌ها، اختصارات و اعداد منبع حفظ شده‌اند. فلش‌های دست‌نویس به ترتیب خواندن بازنویسی شده‌اند. خوانش‌های نامطمئن و اصلاحات دست‌نویس مشخص شده‌اند. تصاویر اصلی در پیوست آمده‌اند.

\IntussusceptionPage{۱}{pic01.jpg}
\textbf{سربرگ:} موضوع: اینتوساسپشن در کودکان؛ تاریخ: / /؛ شمارهٔ برگه: ۱.
\begin{itemize}
\item شایع‌ترین اورژانس شکم در کودکان کم‌سن است، خصوصاً زیر ۲ سال.
\item ۷۵\% موارد ایدیوپاتیک هستند و لید پوینت پاتولوژیک تنها در ۱۵\% موارد یافت می‌شود.
\item ریداکشن:
  \begin{itemize}
  \item به صورت تنقیه با باریم، سایر کنتراست‌ها (انما هیدروستاتیک).
  \item هوا یا \lr{CO\textsubscript{2}}: نیوماتیک انما.
  \end{itemize}
\item پیک سنی ۴–۳۶ ماهگی است.
\item شایع‌ترین علت انسداد روده در این گروه سنی.
\item \textbf{حاشیه:} \lr{3:2}؛ \IntussusceptionUnclear{از نظر \lr{M}} شایع‌تر است.
\item شایع‌ترین \lr{lead point}: \lr{Reactive lymphoid hyperplasia}.
\end{itemize}
بخش پروگزیمال وارد بخش دیستال می‌شود.
\begin{itemize}
\item \lr{Intussusceptum}: بخش پروگزیمال.
\item \lr{Intussuscipiens}: بخش دیستال.
\end{itemize}
\subsection*{تقسیم‌بندی بر اساس بخش}
\begin{itemize}
\item \lr{Ileocolic}: در بچه، ۹۰\% مواقع.
\item \lr{Ileo-ileal, Ileo-ileo-colic, Jejeno-jejunal, \ldots}
\end{itemize}

\IntussusceptionPage{۲}{pic02.jpg}
\textbf{سربرگ:} موضوع: اینتوساسپشن؛ تاریخ: / /؛ شمارهٔ برگه: ۲.

زمانی می‌گوییم ایدیوپاتیک است که در اینتوساسپتوم هیچ \lr{mass}ی یافت نشود.

اگر یک علت احتمالی مثل عفونت وایرال وجود داشته باشد، همچنان ایدیوپاتیک است.

\subsection*{پاتوژنز}
مزانتر به داخل روده کشیده می‌شود
\IntussusceptionArrow احتقان شریان و ورید رخ می‌دهد
\IntussusceptionArrow ادم روده
\IntussusceptionArrow ایسکمی
\IntussusceptionArrow پرفوریشن.

\subsection*{اثر عفونت ویروسی}
\begin{itemize}
\item اینتوساسپشن پترن فصلی دارد و با \lr{AGE} وایرال مرتبط است.
\item با واکسن روتاویروس مرتبط است.
\item ۳۰\% بیماران شرح حال عفونت ویروسی (\lr{URI, GE, \ldots}) دارند.
\item ارتباط قوی با آدنوویروس دارد.
\item می‌تواند بافت لنفاتیک روده را تحریک کند
\IntussusceptionArrow هایپرتروفی \lr{Peyer Patch}.
\item به همین دلیل درمان با \lr{C/S} توصیه می‌شده در مرحلهٔ ریکاورمنت
\IntussusceptionArrow توصیهٔ بالینی نیست (انجام نمی‌دهیم).
\end{itemize}

عفونت‌های باکتریال روده‌ای نیز با اینتوساسپشن همراه هستند.
\begin{itemize}
\item سالمونلا، \lr{E. Coli}، شیگلا، کمپیلوباکتر.
\item ظرف ۱ ماه از انتریت باکتریال.
\end{itemize}

\IntussusceptionPage{۳}{pic03.jpg}
\textbf{سربرگ:} موضوع: اینتوساسپشن؛ تاریخ: / /؛ شمارهٔ برگه: ۳.
\subsection*{لید پوینت‌های شایع}
دایورتیکول‌ها، پولیپ‌ها، تومور، هماتوم، مالفورمیشن‌های عروقی.

\textbf{اصلاح دست‌نویس:} «هماتوم» بالای واژهٔ خط‌خوردهٔ «هامارتوم» نوشته شده است.

اینتوساسپشن در برخی بیماری‌های زمینه‌ای شایع‌تر است:
\begin{itemize}
\item \lr{HSP} (هنوخ)، \lr{IgA Vasculitis}.
\item انگل‌ها (آسکاریس).
\item \lr{CF}.
\end{itemize}
شایع‌ترین: دایورتیکول مکل، سپس پولیپ‌ها.

\subsection*{تظاهرات بالینی تیپیکال}
\begin{itemize}
\item شروع ناگهانی درد شدید، اینترمیتنت، کرامپی، تشدیدشونده.
\item رنگ‌پریدگی.
\item گریه، کنترل نمی‌شود.
\item پاها را به سمت داخل شکم جمع می‌کند.
\item بین اپیزودها حدود ۱۵–۲۰ دقیقه فاصله هست.
\item استفراغ معمولاً به دنبال اولین اپیزود درد شروع می‌شود.
\item در ابتدا غیر صفراوی؛ به دنبال انسداد \IntussusceptionArrow صفراوی.
\end{itemize}
در سمت راست شکم یک تودهٔ سوسیسی لمس می‌شود.

در کمتر از ۲۵\% موارد مدفوع خونی است. در \lr{DRE} با احتمال بیشتری مثبت می‌شود.

\IntussusceptionPage{۴}{pic04.jpg}
\textbf{سربرگ:} موضوع: اینتوساسپشن؛ تاریخ: / /؛ شمارهٔ برگه: \IntussusceptionUnclear{۳؛ شمارهٔ دست‌نویس شبیه ۳ است}.
\begin{itemize}
\item در اکثریتی از بیماران مدفوع \lr{Currant Jelly} دیده می‌شود.
\item ترکیب مخاط و خون.
\item یک یافتهٔ \lr{late} است.
\end{itemize}
وجود خون به صورت \lr{gross} مطرح‌کنندهٔ بیماری روده است.

اگر تصویربرداری \lr{N/C} باشد، به تنهایی کنترااندیکاسیون درمان غیرجراحی نیست.

در سیر بیماری ممکن است بیمار لتارژیک شود و با مننژیت اشتباه گرفته شود.

\subsection*{پرزنتیشن آتیپیک}
تریاد بیماری تنها در ۱۵\% موارد دیده می‌شود:
\begin{enumerate}
\item درد شکم.
\item لمس توده.
\item مدفوع \IntussusceptionUnclear{واژهٔ میانی بین «مدفوع» و «خون» ناخواناست}.
\end{enumerate}
\begin{itemize}
\item ۲۰\% شیرخواران کوچک هیچ درد واضحی ندارند.
\item ۳۰\% بیماران دفع خون و مخاط ندارند.
\item بیماران با سن بیشتر معمولاً فقط درد شکم دارند.
\end{itemize}
در شیرخوار با لتارژی تنها، در کنار \lr{Sepsis} به اینتوساسپشن نیز توجه می‌کنیم.

\IntussusceptionPage{۵}{pic05.jpg}
\textbf{سربرگ:} موضوع: اینتوساسپشن؛ تاریخ: / /؛ شمارهٔ برگه: ۵.

اینتوساسپشن در رودهٔ باریک ممکن است به صورت اتفاقی یافت شود.

حین تصویربرداری‌ها که اغلب موارد گذرا هستند و اهمیت بالینی ندارند.
\begin{itemize}
\item شک از نظر انسداد \IntussusceptionArrow تصویربرداری \lr{Follow up}.
\end{itemize}
\subsection*{در معاینه}
بیرون‌زدگی از مقعد: \lr{Eversion of bowel}.
\begin{latin}
Right lower Quadrant \(\longrightarrow\) Scaphoid (Empty) \(\longrightarrow\) Dance Sign

Right middle/upper \(\longrightarrow\) Sausage, shaped
\end{latin}

\subsection*{تشخیص‌های افتراقی}
\begin{itemize}
\item رکتال بلیدینگ: دایورتیکول مکل، کولیت باکتری، آسیب، ملروتیشن، ولولوس.
\item درد شکم کرامپی با شروع حاد: \lr{GE}، آپاندیسیت، ایسکمی مزانتر، تورشن تخمدان، هرنی اینکارسره، پریتونیت.
\item لتارژی: تروما، سپسیس، اختلال متابولیک، مسمومیت، تشنج.
\end{itemize}

\IntussusceptionPage{۶}{pic06.jpg}
\textbf{سربرگ:} موضوع: اینتوساسپشن؛ تاریخ: / /؛ شمارهٔ برگه: ۶.
\subsection*{منج}
تأیید تشخیص \IntussusceptionArrow مستقیم درمان غیرتهاجمی.

تأیید تشخیص: سونوگرافی / گرافی؛ سونوگرافی انتخابی است.

گرافی: ایستاده، خوابیده؛ برای رد سایر استفاده می‌شود.

نیاز به \lr{ABX} پروفیلاکسی وجود ندارد، مگر ناپایداری همودینامیک یا بیمار \lr{critical ill}.
\begin{itemize}
\item منج: جراحی.
\item غیرجراحی: اغلب بیماران.
\end{itemize}
تکنیک ریداکشن نیوماتیک نسبت به هیدروستاتیک برتر است اگر فلوروسکوپی انجام دهیم.

نیوماتیک + سونو قابل انجام نیست.

\subsection*{منج پس از جااندازی}
\begin{itemize}
\item در یک جااندازی غیرجراحی موفق، نیاز نیست \IntussusceptionArrow \lr{No ABX}.
\item اگر سیر بهترشونده داشت، اگر شک به سپسیس \IntussusceptionArrow \lr{ABX}.
\item اگر شک به پرفوریشن \IntussusceptionArrow \lr{ABX}.
\end{itemize}

\IntussusceptionPage{۷}{pic07.jpg}
\textbf{سربرگ:} موضوع: اینتوساسپشن؛ تاریخ: / /؛ شمارهٔ برگه: ۷.

برخی بیماران پس از جااندازی اندکی تب می‌کنند که قابل انتظار است و نیازمند \lr{ABX} نمی‌باشند.

\subsection*{\lr{Diet}}
بلافاصله بعد از جااندازی مایعات شفاف و در ادامه بر اساس تحمل \lr{PO} باشد.

\lr{NGT} اگر استفراغ‌های مکرر یا شک به انسداد.

تا ۳ ساعت پس از جااندازی اگر \lr{well-appearing} و \lr{PO} را تحمل کند
\IntussusceptionArrow ترخیص.

یک رویکرد دیگر: \lr{observe} برای ۱۲–۲۴ ساعت.

\subsection*{عود}
\begin{itemize}
\item ۴\% ظرف ۴۸ ساعت عود می‌کنند.
\item ۱۰\% جااندازی‌های موفق.
\end{itemize}
بعد از جااندازی نیاز به تصویربرداری نیست مگر علت دیگری.

با عود مثل اپیزود اول برخورد می‌کنیم.

\clearpage
\section*{پیوست: تصاویر اصلی}
\clearpage
\subsection*{تصویر 1}
\begin{center}
\includegraphics[width=\linewidth,height=0.86\textheight,keepaspectratio]{\IntussusceptionImagePath pic01.jpg}
\end{center}
\clearpage
\subsection*{تصویر 2}
\begin{center}
\includegraphics[width=\linewidth,height=0.86\textheight,keepaspectratio]{\IntussusceptionImagePath pic02.jpg}
\end{center}
\clearpage
\subsection*{تصویر 3}
\begin{center}
\includegraphics[width=\linewidth,height=0.86\textheight,keepaspectratio]{\IntussusceptionImagePath pic03.jpg}
\end{center}
\clearpage
\subsection*{تصویر 4}
\begin{center}
\includegraphics[width=\linewidth,height=0.86\textheight,keepaspectratio]{\IntussusceptionImagePath pic04.jpg}
\end{center}
\clearpage
\subsection*{تصویر 5}
\begin{center}
\includegraphics[width=\linewidth,height=0.86\textheight,keepaspectratio]{\IntussusceptionImagePath pic05.jpg}
\end{center}
\clearpage
\subsection*{تصویر 6}
\begin{center}
\includegraphics[width=\linewidth,height=0.86\textheight,keepaspectratio]{\IntussusceptionImagePath pic06.jpg}
\end{center}
\clearpage
\subsection*{تصویر 7}
\begin{center}
\includegraphics[width=\linewidth,height=0.86\textheight,keepaspectratio]{\IntussusceptionImagePath pic07.jpg}
\end{center}
\IntussusceptionFinish
